% ORACLE C RESULT.
% This fragment addresses only the algebraic marked-coordinate prerequisite.
% It proves no dispersion estimate and makes no claim toward
% eq:oracleA-sdc or eq:oracleB-mh2.

\subsection{A fixed toric marked coordinate}
\label{ssec:oracleC-result}

Put
\begin{equation}
 \Lambda(x,y,z)=
 \frac{(1+x)(1+y)(1+z)
  \bigl((1+y)(1+z)+xyz\bigr)}{xyz}
 \in\mathbb Z[x^{\pm1},y^{\pm1},z^{\pm1}].
 \label{eq:oracleC-Lambda}
\end{equation}
This one Laurent polynomial gives the required marked scalar.  In
particular, it avoids both loci ruled out by
Propositions~\ref{prop:oracleB-two-zero-loci}
and~\ref{prop:oracleB-complexity}: it is neither an evaluation of $H_p$
nor an evaluation of its factor $B_p$.

\begin{lemma}[Fixed constant-term realization]
\label{lem:oracleC-constant-term}
For every integer $n\geq0$,
\begin{equation}
 \operatorname {CT}_{x,y,z}\Lambda(x,y,z)^n
 =\sum_{r=0}^n\binom nr^2\binom{n+r}{r}^2=b_n.
 \label{eq:oracleC-constant-term}
\end{equation}
Moreover, every exponent of each of $x,y,z$ occurring in~$\Lambda$ lies
between $-1$ and $1$.
\end{lemma}

\begin{proof}
Write
\[
 A=(1+x)(1+y)(1+z),\qquad
 E=(1+y)(1+z)+xyz.
\]
Then the constant term in question is
\[
 [x^{n}y^{n}z^{n}]A^nE^n.
\]
If exactly $k$ of the $n$ factors selected in $E^n$ contribute $xyz$,
the $x$-coefficient contributes a second factor $\binom nk$, while the
$y$- and $z$-coefficients both contribute
$\binom{2n-k}{n-k}$.  Hence
\[
 \operatorname {CT}\Lambda^n
 =\sum_{k=0}^n\binom nk^2\binom{2n-k}{n-k}^2.
\]
The substitution $r=n-k$ gives~\eqref{eq:oracleC-constant-term}.
The exponent assertion follows directly from the numerator of
\eqref{eq:oracleC-Lambda}, whose degree in each variable is between $0$
and $2$, followed by division by $xyz$.
\end{proof}

\begin{theorem}[The marked coordinate]
\label{thm:oracleC-marked-coordinate}
Let $p\geq5$ be prime.  For $0\leq j\leq p-2$, define the scalar
\begin{equation}
 c_{p,j}^{\rm tor}
 :=-\sum_{(x,y,z)\in(\mathbb F_p^\times)^3}
       \Lambda(x,y,z)^j\in\mathbb F_p,
 \label{eq:oracleC-toric-coordinate}
\end{equation}
where the zeroth power is the constant Laurent polynomial $1$, including
at the zeros of $\Lambda$.  Define separately
\[
 c_{p,p-1}^{\rm tor}:=1.
\]
Then, for every $0\leq j\leq p-1$,
\begin{equation}
 c_{p,j}^{\rm tor}=b_j\pmod p.
 \label{eq:oracleC-coordinate-equality}
\end{equation}
Consequently
\begin{equation}
 c_{p,j}^{\rm tor}=0
 \quad\Longleftrightarrow\quad
 b_j\equiv0\pmod p.
 \label{eq:oracleC-marked-zero}
\end{equation}
\end{theorem}

\begin{proof}
For an integer $e$,
\[
 \sum_{a\in\mathbb F_p^\times}a^e
 =\begin{cases}-1,&p-1\mid e,\\0,&p-1\nmid e,
 \end{cases}
 \qquad\text{in }\mathbb F_p.
\]
By Lemma~\ref{lem:oracleC-constant-term}, every exponent in
$\Lambda^j$ belongs to $[-j,j]$ in each coordinate.  If
$0\leq j\leq p-2$, the only multiple of $p-1$ in this interval is zero.
Three applications of torus orthogonality therefore give
\[
 \sum_{(\mathbb F_p^\times)^3}\Lambda^j
 =(-1)^3\operatorname {CT}\Lambda^j=-b_j,
\]
which proves~\eqref{eq:oracleC-coordinate-equality} on this range.  At the
remaining endpoint, the reflection in Remark~\ref{rem:orbit} gives
$b_{p-1}=b_0=1$ modulo $p$, exactly the separately defined value.
\end{proof}

\begin{remark}[The endpoint is not a cosmetic convention]
\label{rem:oracleC-endpoint}
The raw power sum at $j=p-1$ cannot replace the separate definition.
If
\[
 T=(\mathbb G_m)^3,\qquad U=T\setminus V(\Lambda),
\]
then direct inclusion--exclusion gives
\begin{equation}
 \#U(\mathbb F_p)=p^3-7p^2+17p-14.
 \label{eq:oracleC-U-count}
\end{equation}
Indeed, the union of $x=-1$, $y=-1$, and $z=-1$ has
$(p-1)^3-(p-2)^3$ points.  Off $y=-1$ and $z=-1$, the fourth zero
equation determines a unique nonzero $x$ and has $(p-2)^2$ points; its
intersection with the first union has $p-2$ points, given by
$x=-1$ and $1+y+z=0$.  Subtracting this zero locus from $(p-1)^3$
gives~\eqref{eq:oracleC-U-count}.
Thus
$-\sum_{T(\mathbb F_p)}\Lambda^{p-1}=-\#U(\mathbb F_p)\equiv14\pmod p$,
not $b_{p-1}=1$ in general.  At $p=7$ the raw endpoint is even zero.
For $j=0$, by contrast, the convention in
\eqref{eq:oracleC-toric-coordinate} gives
$-(p-1)^3=1$ in $\mathbb F_p$.
\end{remark}

\subsubsection{Uniform geometric complexity}

For $1\leq j\leq p-2$, let $E_p=\mathbb Q(\mu_{p-1})$.  Choose a prime
$\mathfrak p\mid p$ and a Teichm\"uller character
$\omega_p\colon\mathbb F_p^\times\to\mu_{p-1}$ such that
$\omega_p(a)\equiv a\pmod{\mathfrak p}$, and put
$\chi_{p,j}=\omega_p^j$.  Fix $\ell=2$ (recall that $p\geq5$), choose an
embedding $E_p\hookrightarrow\overline{\mathbb Q}_2$, and let
$\mathcal K_{\chi_{p,j}}$ be the associated rank-one Kummer sheaf, with
the Frobenius convention chosen to have trace $\chi_{p,j}$.  On the fixed
threefold $U$ put
\begin{equation}
 \mathcal L_{p,j}=\Lambda^*\mathcal K_{\chi_{p,j}}.
 \label{eq:oracleC-Kummer-object}
\end{equation}
The Grothendieck trace formula gives the cyclotomic integer
\begin{equation}
 \widetilde c_{p,j}^{\rm tor}
 =-\operatorname {Tr}_{\rm alt}\!\left(
   \operatorname {Frob}_p\mid
   R\Gamma_c(U_{\overline{\mathbb F}_p},\mathcal L_{p,j})
  \right)
 =-\sum_{u\in U(\mathbb F_p)}\chi_{p,j}(\Lambda(u)).
 \label{eq:oracleC-Frobenius-coordinate}
\end{equation}
Its reduction at the chosen prime above $p$ is
$c_{p,j}^{\rm tor}=b_j$.  Both endpoints may instead be represented by
the constant sheaf on $T$: its negative Frobenius supertrace is
$-(p-1)^3$, whose reduction is $1$.  Thus their cohomological complexity
is fixed as well.  Using the trivial Kummer sheaf extended by zero from
$U$ would give the wrong value in Remark~\ref{rem:oracleC-endpoint}.

\begin{proposition}[Uniform complexity of the toric object]
\label{prop:oracleC-uniform-complexity}
The ambient variety and rational map in
\eqref{eq:oracleC-Kummer-object} are independent of $p$ and $j$.  The
coefficient sheaf has rank one.  There is an absolute constant
$B_\Lambda$ such that
\begin{equation}
 \sum_i\dim H_c^i
 \bigl(U_{\overline{\mathbb F}_p},\mathcal L_{p,j}\bigr)
 \leq B_\Lambda
 \label{eq:oracleC-Betti-bound}
\end{equation}
for all $p\geq5$ and $1\leq j\leq p-2$.  In particular, both the total
degree of the corresponding cohomological $L$-factor and the ramification
complexity are $O(1)$ uniformly in $(p,j)$.  This is geometric complexity:
as for the Kummer twists in~\eqref{eq:oracleB-crystalline-mellin}, it does
not count the degree of the varying cyclotomic coefficient field $E_p$.
\end{proposition}

\begin{proof}
In $(\mathbb P^1)^3$, the principal divisor of $\Lambda$ has the four
simple zero components
\[
 x=-1,\qquad y=-1,\qquad z=-1,\qquad
 \overline Q:\
 X_0(Y_0+Y_1)(Z_0+Z_1)+X_1Y_1Z_1=0
\]
and the six simple pole components
\[
 x=0,\ x=\infty,\ y=0,\ y=\infty,\ z=0,\ z=\infty.
\]
The fourth zero component is smooth and geometrically irreducible.  Indeed,
its two coefficients as a linear form in $(X_0,X_1)$ are coprime, and a
direct check of the six displayed partial derivatives has no projective
common zero.  On the torus, smoothness is already immediate from
$\partial((1+y)(1+z)+xyz)/\partial x=yz$.
Thus the input divisor has ten fixed components of fixed multidegrees.
The full divisor is not asserted to be simple normal crossings.

More explicitly, the change of coordinates
\[
 u=\frac{1+y}{y},\qquad v=\frac{1+z}{z},\qquad
 w=-\frac{x}{uv}
\]
identifies $U$ with the fixed toric-arrangement complement
\[
 (\mathbb G_m)^3\setminus
 \{u=1,\ v=1,\ w=1,\ uvw=1\}.
\]
In particular, a finite stratification depends only on this four-member
arrangement, not on the character.

Every $\chi_{p,j}$ has order dividing $p-1$.  Its pullback is therefore
tame along every divisorial valuation and has Swan conductor zero.  Take
one log resolution of this fixed integral divisor and spread it out away
from finitely many characteristics.  The resulting stratification has a
fixed finite number of strata.  At every good prime, tame specialization
for this log-smooth pair applies because the local monodromy has order
prime to $p$; it compares the relevant rank-one representation with one on
the characteristic-zero fiber.  Fix a finite CW model of that complex
fiber.  Its cellular cochain groups have dimension at most the number of
cells for every rank-one representation, independently of its monodromy,
and hence give~\eqref{eq:oracleC-Betti-bound}.
The finitely many excluded characteristics contribute only finitely many
additional $(p,j)$ and may be absorbed into $B_\Lambda$.  This proves a
uniform bound without claiming that the original ten-component divisor is
normal crossings or that the cohomology is concentrated in one degree.
\end{proof}

\begin{remark}[What vanishes]
\label{rem:oracleC-integral-warning}
Equation~\eqref{eq:oracleC-marked-zero} concerns the reduction of the
cyclotomic trace~\eqref{eq:oracleC-Frobenius-coordinate} modulo a prime
above $p$.  It does not assert that the characteristic-zero Frobenius
trace itself vanishes.  For example, at $(p,j)=(11,5)$ the quadratic
character sum gives $\widetilde c_{11,5}^{\rm tor}=-33\ne0$, while its
reduction and $b_5$ both vanish modulo $11$.  This integral distinction is
essential in any later attempt to prove~\eqref{eq:oracleB-mh2}.
\end{remark}

\subsubsection{Audits of the two originally suggested coordinates}

The toric construction above is not a disguised bounded-length transfer
word.  With $B_n=(n!)^3b_n$ and the matrix $M(n)$ from
Remark~\ref{rem:orbit}, put $T_0=I$ and
$T_j=M(j-1)\cdots M(0)$.  The recurrence gives the exact identity
\begin{equation}
 T_j=\begin{pmatrix}B_j&0\\B_{j-1}&0\end{pmatrix}
 \qquad(j\geq1).
 \label{eq:oracleC-transfer-marker}
\end{equation}
Thus the marked entry is $(1,1)$, not $(1,2)$, and it detects $b_j=0$
for $j<p$.  Its matrix size is fixed, but its word length is $j$.  Exact
symbolic multiplication of
$R_\ell(x)=M(x+\ell-1)\cdots M(x)$ gives, for $\ell\geq2$,
\[
 \deg R_\ell=
 \begin{pmatrix}3\ell&3\ell+3\\3\ell-3&3\ell\end{pmatrix}.
\]
For completeness, if $A_\ell=(R_\ell)_{11}$ and $\lambda_\ell$ is its
leading coefficient, then
\[
 A_{\ell+1}=P(x+\ell)A_\ell-(x+\ell)^6A_{\ell-1},
 \qquad
 \lambda_{\ell+1}=34\lambda_\ell-\lambda_{\ell-1},
 \quad(\lambda_0,\lambda_1)=(1,34).
\]
The $\lambda_\ell$ are positive and strictly increasing, so the two
leading terms never cancel.  The remaining entries satisfy
\[
 (R_\ell)_{21}=A_{\ell-1}(x),\qquad
 (R_\ell)_{12}=-x^6A_{\ell-1}(x+1),\qquad
 (R_\ell)_{22}=-x^6A_{\ell-2}(x+1),
\]
for $\ell\geq2$.  This proves the displayed degree matrix rather than
extrapolating it from finitely many products.
Likewise, the standard factorial gauge has marked coefficient
\[
 \frac{N_h(m)}{\prod_{s=2}^h(m+s)^3}.
\]
The endpoint identity used in Remark~\ref{rem:content},
\[
 N_h(-s)=(-1)^{s-1}B_{s-1}B_{h-s}\ne0
 \qquad(1\leq s\leq h),
\]
shows that it has exactly $h-1$ uncancelled finite poles, each of order
three.  These facts rule out the direct word and this direct gauge; they
are not a no-go theorem for all constructions, as
Theorem~\ref{thm:oracleC-marked-coordinate} demonstrates.

The raw Mellin coordinate from
Proposition~\ref{prop:oracleB-two-zero-loci} has a different restricted
obstruction.  If $g$ is a primitive root, its corrected group-algebra
element is
\begin{equation}
 Q_p(X)=-\sum_{r=0}^{p-2}H_p(g^r)X^{-r}
 \quad\text{in }\mathbb F_p[X^{\pm1}]/(X^{p-1}-1).
 \label{eq:oracleC-group-algebra}
\end{equation}
For $1\leq j\leq p-2$, $Q_p(g^j)=b_j$.  The factorization in
Remark~\ref{rem:squareness} implies, without any squarefreeness or
coprimality assumption,
\[
 \#\{t\in\mathbb F_p^\times:H_p(t)=0\}
 \leq\deg B_p+2\varepsilon_p
 =\frac{p-1}{2}+\varepsilon_p,
\]
and hence
\[
 \#\operatorname {supp}Q_p
 =\#\{t\in\mathbb F_p^\times:H_p(t)\ne0\}
 \geq\frac{p-1}{2}-\varepsilon_p
 \geq\frac{p-3}{2}.
\]
Hence the raw periodic coordinate has growing group-algebra support and
growing minimal cyclic constant-coefficient recurrence order: the distinct
characters in~\eqref{eq:oracleC-group-algebra} are linearly independent
because the $(p-1)$-point Fourier/Vandermonde matrix is invertible in
$\mathbb F_p$.  This is a
no-go only for exact raw-value realizations of that form, not for the
zero-equivalent toric coordinate.

Finally, bounded complexity of a Kummer factor on the original $t$-line
does not by itself repair~\eqref{eq:oracleB-crystalline-mellin}.  At
$p=31$, the two roots of the conifold discriminant are $14$ and $20$,
with $H_{31}(14)=H_{31}(20)=7$ and $b_8=b_{22}=0$.  For $j=8$ the omitted
bad-fiber contribution is
\[
 7(14^{-8}+20^{-8})=11\pmod {31},
\]
so the Mellin sum restricted to the smooth locus is $20$, not zero; the
same failure occurs for $j=22$.  Boundary terms require integral control.

\subsubsection{Quantifier boundary}

The family~\eqref{eq:oracleC-Kummer-object} is a uniformly bounded family
of cohomological spaces indexed by the tame character $j$.  It is not a
single sheaf on a ``$j$-line'' whose trace function is
$j\mapsto b_j$.  Therefore a number of $j$-singularities or a conductor
of such a $j$-line sheaf is not defined.  In particular,
Proposition~\ref{prop:oracleC-uniform-complexity} does not turn the
moving-character problem into an ordinary FKM bilinear estimate.
It is also an alternative toric Kummer realization of the marked scalar,
not a construction of, or a comparison with, the transcendental crystal
$\mathcal V_{\mathrm{tr},p}$ in
\eqref{eq:oracleB-crystalline-mellin}.  Thus it supplies exactly the
algebraic coordinate requested here, not the rest of that crystalline
Mellin package.

\paragraph{STALL REPORT.}
The literal algebraic goal G1 is achieved by
\eqref{eq:oracleC-toric-coordinate}: the marked scalar is exactly $b_j$
modulo $p$, and its ambient dimension, coefficient rank, divisor data,
Swan conductors, and total cohomological dimension are uniformly bounded
in $(p,j)$.  G2 is false in its universal form and is not claimed; the
transfer and raw-Mellin statements above are restricted no-go results.
G3 is deliberately not attempted.  No estimate for
\eqref{eq:oracleA-sdc}, no short-arc bound, no two-characteristic
dispersion estimate~\eqref{eq:oracleB-mh2}, and no reduction of any of
them to an existing theorem is proved here.  The remaining obstruction is
precisely horizontal control of the moving characters and their reductions,
not the existence of an algebraic marked coordinate.
